\section{Discussion}\label{sec:discussion}

Context types provide a type-based theoretical foundation for safety and reachability verification; however, there is an intrinsic difficulty in introducing a decidable bidirectional typing algorithm like those for standard refinement types. The core challenges come from the following three aspects:
\begin{itemize}
    \item Not totally term-syntax guided. Bidirectional typing algorithms and other typing algorithms typically leverage the syntax of the program to guide typing; however, within a bunched typing setting, the same term syntax admits multiple different typing rules (e.g., \textsc{T-Let} and \textsc{T-LetD}). This forces each typing algorithm to search through different choices, causing the case analysis to explode. Allowing both entangled and independent ways to construct the type context also yields a tree-like context structure. Like bunched implication (e.g., separation logic), this forces the algorithm to construct and decompose the context structure in a large search space.
    \item Not totally type-syntax guided. The syntax of the type we type-check against also cannot provide enough guidance. Since function types can have different modalities for parameters and returns, the polarity of demonic and angelic modalities can alternate during the type-checking process.
    For example, consider the following example, which mimics using a demonic sum type to express a reachability verification task introduced in \autoref{sec:typing}:
    \begin{align*}
       &\Gamma \ctcomma f{:}(\urt{\Nat}{\vnu > 0}\ctsarr\ort{\Nat}{\vnu < 0} \ctoplus \ort{\Nat}{\top}) \vdash
       \\&\quad \zlet{x}{e_x}{f\;x} : \ort{\Nat}{\vnu > 0} \ctoplus \ort{\Nat}{\top}
    \end{align*}\noindent
    Although the right-hand-side type is demonic, we should type-check the term $e_x$ against the angelic type $\urt{\Nat}{\vnu > 0}$. 
    Alternatively, when function $f$ has demonic type, the term $e_x$ should be typed against the demonic type $\ort{\Nat}{\vnu > 0}$:
    \begin{align*}
        &\Gamma \ctcomma f{:}(\ort{\Nat}{\vnu > 0}\ctsarr\ort{\Nat}{\vnu < 0} \ctoplus \ort{\Nat}{\top}) \vdash
        \\&\quad \zlet{x}{e_x}{f\;x} : \ort{\Nat}{\vnu > 0} \ctoplus \ort{\Nat}{\top}
    \end{align*}\noindent
    In fact, whether the binding expression $e_x$ should be typed demonic or angelic depends on the body expression in the let-binding rather than on the right-hand-side type. This makes building a decidable typing algorithm require analyzing the continuation, which can be arbitrarily complex.
    \item Undecidable subtyping checking. Our refinement type system allows arbitrary mixing of demonic and angelic modalities, and our semantic subtyping rules require checking verification conditions that may have arbitrary quantifier alternation, which is undecidable in general.
\end{itemize}

We would like to highlight that a decidable bidirectional typing algorithm is still possible to achieve with a limited version of context types. The coverage type system~\cite{CoverageType} provides an example, where all function types are persistent with demonic parameter types and angelic return types. These constraints solve the three issues by (1) using variable references in type qualifiers to express entanglement, i.e., the dependent context $x{:}\urt{b}{\top} \ctcomma y{:}\urt{b}{\top}$ is illegal in coverage type system since their qualifiers do not refer to each other; (2) only allowing the right-hand-side type to be angelic, and providing specialized application rules for angelic arguments and function parameter types; (3) requiring that local angelic variables appear only after demonic parameters in the type context, which leads verification conditions to fall in EPR (i.e., in a $\forall^*\exists^*$ form), which is decidable.

\section{Related Work}\label{sec:related}

\paragraph{Bunched typing and bunched implication.}
Although Bunched implication is wilded used in program logics, such as seperation logic~\cite{Iris}, the corresponding typing discipline has rarely been explored. Besides the original work by O'Hearn~\cite{BunchedType}, our work is the only bunched type system we are aware of. Context adapts the bunched typing into refinement type system, and equip it with richer modalities (i.e., demomic and angelic modalities, sum type, and persistence modality) to unify safety and reachability verification.

\paragraph{Reachability verification.}
Recent work on reachability verification in type systems includes
Ramsay's ill-types~\cite{Ill-types,RamsayIncorType}, which refute
typing formulas to show that ill-formed programs do not evaluate. Reachability types~\cite{ReachabilityTypes,ReachabilityTypesBidirectional,ReachabilityTypesTypestate,ReachabilityTypesLifetimes,ReachabilityTypesCyclic,ReachabilityTypesLogicalRelations,ReachabilityTypesPolymorphic} form a series of works that track aliasing of mutable references in higher-order functional programs by annotating each type with a reachability set of variable names that the reference may alias.
Coverage types~\cite{CoverageType} underapproximates
test-input-generator behavior to establish coverage completeness.
All of these works are more algrithmical and not provides a thoeritical foundation such as an underline semanticaly denotation as our capability algebra and context Kirpke logic in our appoarch. As we discussed in \autoref{sec:discussion}, coverage type can be treated as a limited version of context types that to support decidable bidirectional typing algorithm.

In program logics, the literature is larger and has focused primarily
on \emph{incorrectness reasoning}.
Underapproximate logics~\cite{OHP19,RBD+20,ReverseHoareLogic,
HyperHoareLogic,QuantitativeWeakestHyperPre,UTurn,BackwardIncorrectness}
study judgments that guarantee the existence of executions rather than
universal safety.
However, all of these program logic has a unified pre- and post-condition in a single modality (angelic or demonic), which doesn't support construct or decompose the context structure as in our work.
Outcome logic~\cite{ZDS23,Zil25TOPLAS} is the closest to our work
because it also unifies safety and reachability through semiring-weighted
outcomes and branch composition $\varphi \oplus \psi$.
Like our $\tau_1 \ctoplus \tau_2$, its $\oplus$ adjoins control-flow
alternatives, but outcome logic's base semantics is demonic and
recovers incorrectness via $\langle \chi \oplus \top \rangle$ rather
than through native underapproximate modalities $\ort{b}{\phi}$ and
$\urt{b}{\phi}$ on the same refinement.
These logics are command-first and store-based, whereas we give a
complementary higher-order account through context types and bunched
witness structure.

\paragraph{May- and must-style program analysis.} The program analysis works validate the application domain of coontext types.
Compositional may- and must-style program analysis~\cite{MayMust}
alternates overapproximate may-summaries with underapproximate
must-summaries so that callee information can inform caller analysis.
Refutation is a type-level specialization of this pattern, since a failed
refinement type check does not by itself establish that a bug is
reachable.
Refinement type refutations~\cite{RefinementTypeRefutations} address
this gap by deriving \emph{must}-style witnesses that explain why
typing failed, often through concrete must-instantiations of subterms,
while symbolic-execution tooling~\cite{RefinementTypeSymbolicExecution}
searches for such counterexamples externally.
Again, our context types provides a unified theoretical foundation for these algorithms which revists the duality between safety and reachability reaosning.

\paragraph{Context-aware typing.}
Coeffect and graded modal type systems~\cite{Coeffects,
Katsumata+14, GaboardiCombining, OrchardQuantitative} annotate
programs with semiring-indexed context information such as usage,
liveness, and security levels, interpreted via graded (co)monads.
Like our context types, they treat contextual information as a typed
resource with a compositional algebra, but they track how a
computation \emph{demands} from its environment rather than how a
typing context \emph{interprets} value sets under demonic or angelic
modalities.
They therefore do not internalise safety/reachability duality on the
same base refinement, nor the bunched capability structure needed for
higher-order witness independence (\autoref{subsec:disprove}).